\begin{lemma}
\label{editorial-source-lemma-geometrically-irreducible-any-base-change}
Let $k$ be a field.
Let $S$ be a geometrically irreducible $k$-algebra.
Let $R$ be any $k$-algebra.
The map
$$
\Spec(R \otimes_k S) \longrightarrow \Spec(R)
$$
induces a bijection on irreducible components.
\end{lemma}

\begin{proof}
Write $A=R\otimes_k S$ and retain
$f:\Spec(A)\to\Spec(R)$. The
original $S$ is nonzero, since
geometric irreducibility includes
nonemptiness. If $R=0$, then
$A=0$ and both sets of components
are empty, so the assertion
holds. Otherwise, for every
$\mathfrak p\in\Spec(R)$ the
original quotient and fibre maps
give
$$
A/\mathfrak pA\simeq(R/\mathfrak p)\otimes_k S
\longrightarrow\kappa(\mathfrak p)\otimes_k S.
$$
The quotient isomorphism sends
$[r\otimes s]$ to $\overline r\otimes s$
and has inverse
$\overline r\otimes s\mapsto[r\otimes s]$.
The arrow is injective, since
$R/\mathfrak p\to\kappa(\mathfrak p)$
is an injection of $k$-vector
spaces and tensoring by $S$
preserves it. Its target has
prime nilradical by the given
geometric irreducibility of $S$.
The inverse image of this
nilradical is precisely the
nilradical of the source, by
injectivity, and is therefore
a proper prime. In particular
$A/\mathfrak pA$ has irreducible
nonempty spectrum.

Consequently the original ideal
$$
\sigma(\mathfrak p)=\sqrt{\mathfrak pA}
$$
is prime, and $f^{-1}(V(\mathfrak p))$
is the irreducible closed subset
$V(\sigma(\mathfrak p))$.
Its contraction is exactly
$\mathfrak p$. To verify this,
the injection
$R/\mathfrak p\to(R/\mathfrak p)\otimes_k S$
can be split as a linear map
by choosing a $k$-linear
functional $S\to k$ taking
$1$ to $1$. Hence an element
of $R/\mathfrak p$ becomes
nilpotent in that tensor
exactly when it was nilpotent
in the original domain, that
is, exactly when it is zero.
This proves the contraction
claim with the given scalar
map.

The map $R\to A$ is flat: a
$k$-basis of $S$ gives a free
$R$-module basis of $A$.
Going down therefore sends
every minimal prime of $A$
to a minimal prime of $R$.
If $\mathfrak p$ is minimal
in $R$, choose a minimal
$\mathfrak q_0\subseteq\sigma(\mathfrak p)$
of $A$. Its contraction is
a minimal prime contained
in $\mathfrak p$, hence is
exactly $\mathfrak p$.
It contains $\mathfrak pA$,
so it contains
$\sqrt{\mathfrak pA}=\sigma(\mathfrak p)$.
Thus equality holds, proving
that $\sigma(\mathfrak p)$
is minimal. Conversely, for
any minimal $\mathfrak q$
of $A$, put
$\mathfrak p=\mathfrak q\cap R$.
Then $\sigma(\mathfrak p)\subseteq\mathfrak q$,
so minimality makes them equal.
These actual inverse maps on
minimal primes prove the
bijection on components via
Lemma \ref{editorial-source-lemma-irreducible}.

The construction gives more:
$\sigma$ is a continuous section
of $f$ on the entire spectra,
selecting the generic point
of each original fibre.
For every $\mathfrak p$, its
prime is contained in every
prime of $A$ over $\mathfrak p$,
and its contraction equals
$\mathfrak p$, so it is the
unique generic fibre prime.
For an original $a\in A$ one
has the exact equality
$$
\sigma^{-1}(D(a))=f(D(a)).
$$
If a prime above $\mathfrak p$
does not contain $a$, the
smaller $\sigma(\mathfrak p)$
does not contain it. Conversely
if $a\notin\sigma(\mathfrak p)$,
this prime itself is a point
of $D(a)$ above $\mathfrak p$.
Lemma \ref{editorial-source-lemma-map-into-tensor-algebra-open}
makes the right side open,
proving continuity. Since
$f\sigma$ is the identity,
the section is a topological
embedding, with inverse the
restriction of $f$ to its
image. That image is dense:
every nonempty basic open
$D(a)$ has nonempty image,
and the displayed equality
puts a section point in it.
Also
$$
\overline{\{\sigma(\mathfrak p)\}}
=V(\sqrt{\mathfrak pA})
=f^{-1}(V(\mathfrak p)).
$$
Restricting the two continuous
inverse maps to the minimal
primes proves a homeomorphism
of their subspaces in the
respective spectra, strengthening
the stated component bijection.

Finally this section is
compatible with arbitrary
original $k$-algebra base change
$\theta:R\to R'$. If $\mathfrak p'$
contracts to $\mathfrak p$,
the induced map of fields
$\kappa(\mathfrak p)\to\kappa(\mathfrak p')$
gives an injection
$$
\kappa(\mathfrak p)\otimes_k S
\longrightarrow\kappa(\mathfrak p')\otimes_k S.
$$
The inverse image of the
nilradical on the right is
the nilradical on the left,
since injective ring maps
reflect nilpotence. The two
quotient-fibre descriptions
above therefore show that
$\sigma_{R'}(\mathfrak p')$
contracts to
$\sigma_R(\mathfrak p)$ in
$R\otimes_k S$. This proves
the stated compatibility
through the original tensor
map $r\otimes s\mapsto\theta(r)\otimes s$.
The section is a continuous
map of these spectra; no
extra residue-field
isomorphisms are asserted.
\end{proof}